%----------------------------------------------------------------------------------------
%	1. GÓI NGÔN NGỮ & FONT
%----------------------------------------------------------------------------------------
\usepackage[T5,T1]{fontenc}
\usepackage{mathptmx}
\usepackage[vietnamese]{babel}
\usepackage{booktabs,tabularx}
\usepackage{algorithm}
%----------------------------------------------------------------------------------------
%	2. CÁC GÓI CƠ BẢN
%----------------------------------------------------------------------------------------
\usepackage{amsmath}
\usepackage{graphicx}
\usepackage{array,booktabs,multirow}
\usepackage{siunitx} % canh lề số đẹp
\sisetup{
  round-mode=places,round-precision=3,
  table-number-alignment = center,
  detect-weight=true}

\usepackage{subcaption}
\usepackage{float}
\usepackage{setspace}
\usepackage{longtable}
\usepackage{tabularray}
\usepackage{tikz}
\usetikzlibrary{calc}
\usepackage{etoolbox}

%----------------------------------------------------------------------------------------
%	3. BỐ CỤC & SIÊU LIÊN KẾT
%----------------------------------------------------------------------------------------
\usepackage{geometry}
\geometry{
  a4paper,
  top=20mm,
  bottom=20mm,
  left=30mm,
  right=20mm,
  headsep=10mm,
  footskip=10mm,
}
\definecolor{cobalt}{rgb}{0.0, 0.28, 0.67}
\definecolor{carmine}{rgb}{0.59, 0.0, 0.09}
\usepackage[unicode, colorlinks=false, pdfborder={0 0 0}]{hyperref}

%----------------------------------------------------------------------------------------
%	4. ĐỊNH DẠNG TIÊU ĐỀ (CHƯƠNG, MỤC, CAPTION)
%----------------------------------------------------------------------------------------
\usepackage{titlesec}
\usepackage{caption}
\setcounter{secnumdepth}{3} % đánh số tới \subsubsection  -> 1.3.1
\setcounter{tocdepth}{3}    % hiển thị tới \subsubsection trong TOC

% Định dạng tiêu đề Chương
\titleformat{name=\chapter}[block]
  {\normalfont\fontsize{14pt}{20pt}\bfseries\centering}
  {\fontsize{13pt}{16pt}\selectfont\MakeUppercase{Chương \thechapter.}}
  {1em}
  {\fontsize{14pt}{20pt}\selectfont\MakeUppercase}
\titleformat{name=\chapter,numberless}[block]
  {\normalfont\fontsize{14pt}{20pt}\bfseries\centering}
  {}{0em}{\MakeUppercase}
\titlespacing*{\chapter}{0pt}{-10pt}{14pt}

% Định dạng tiêu đề Section & Subsection
\titleformat{\section}[hang]
  {\normalfont\fontsize{13pt}{15pt}\selectfont\bfseries}
  {\thesection}{1em}{}
\titleformat{\subsection}[hang]
  {\normalfont\fontsize{13pt}{16.5pt}\selectfont\bfseries}
  {\thesubsection}{1em}{}

% Định dạng tiêu đề cho hình, bảng
\captionsetup{font=normalsize, labelfont=bf, labelsep=colon}

%----------------------------------------------------------------------------------------
%	5. ĐỊNH DẠNG MỤC LỤC & CÁC DANH SÁCH
%----------------------------------------------------------------------------------------
\usepackage{tocloft}
\usepackage{textcase} % \MakeTextUppercase - xử lý đúng chữ có dấu tiếng Việt khi in hoa

% Định dạng tiêu đề của các trang danh sách
\makeatletter
\renewcommand{\@cftmaketoctitle}{%
  \clearpage
  {\centering\bfseries\fontsize{14pt}{20pt}\selectfont \MakeTextUppercase{\contentsname}\par}
  \nobreak\vskip 14pt
  {\normalsize\hfill Trang}\par\vskip 4pt
}
\renewcommand{\@cftmakeloftitle}{%
  \clearpage
  {\centering\bfseries\fontsize{14pt}{20pt}\selectfont \MakeTextUppercase{\listfigurename}\par}
  \nobreak\vskip 14pt
}
\renewcommand{\@cftmakelottitle}{%
  \clearpage
  {\centering\bfseries\fontsize{14pt}{20pt}\selectfont \MakeTextUppercase{\listtablename}\par}
  \nobreak\vskip 14pt
}
\makeatother

% Định dạng các dòng trong Mục lục
\renewcommand{\cftchappresnum}{Chương }
\renewcommand{\cftchapaftersnum}{.}
\setlength{\cftchapnumwidth}{8em}
\renewcommand{\cftsecleader}{\cftdotfill{\cftdotsep}}
\renewcommand{\cftsubsecleader}{\cftdotfill{\cftdotsep}}
\renewcommand{\cftdotsep}{3}

% Mục cấp Chương (Chương 1, Lời cam đoan, Kết luận...) — in hoa, in đậm
% Lưu ý: không dùng \renewcommand{\cftchapfont}{\bfseries\MakeUppercase} vì
% cách này chỉ "ăn" 1 ký tự đầu, làm vỡ dấu tiếng Việt — phải ghi đè \l@chapter
% trực tiếp để bọc đúng dấu ngoặc cho \MakeTextUppercase nhận đủ toàn bộ tên chương.
\renewcommand{\cftchapfont}{\bfseries\fontsize{12pt}{14pt}\selectfont}
\renewcommand{\cftchappagefont}{\bfseries\fontsize{12pt}{14pt}\selectfont}
\renewcommand{\cftchapleader}{\cftdotfill{\cftdotsep}}

\makeatletter
\renewcommand*{\l@chapter}[2]{%
  \ifnum \c@tocdepth >\m@ne
    \addpenalty{-\@highpenalty}%
    \vskip \cftbeforechapskip
    {\leftskip \cftchapindent\relax
     \rightskip \@tocrmarg
     \parfillskip -\rightskip
     \parindent \cftchapindent\relax\@afterindenttrue
     \interlinepenalty\@M
     \leavevmode
     \@tempdima \cftchapnumwidth\relax
     \let\@cftbsnum \cftchappresnum
     \let\@cftasnum \cftchapaftersnum
     \let\@cftasnumb \cftchapaftersnumb
     \advance\leftskip \@tempdima \null\nobreak\hskip -\leftskip
     {\cftchapfont \MakeTextUppercase{#1}}\nobreak
     \cftchapfillnum{#2}}%
  \fi}%
\makeatother

% Mục cấp 1.1 (section) — chỉ in đậm, không in hoa
\renewcommand{\cftsecfont}{\bfseries}
\renewcommand{\cftsecpagefont}{\bfseries}

% Làm đều khoảng cách các dòng trong Mục lục
\setlength{\cftbeforechapskip}{\parskip}
\setlength{\cftbeforesecskip}{\parskip}
\setlength{\cftbeforesubsecskip}{\parskip}

%----------------------------------------------------------------------------------------
%	6. CÁC GÓI MỞ RỘNG (CODE, GIẢI THUẬT)
%----------------------------------------------------------------------------------------
\usepackage{algorithm}
\usepackage{algpseudocode}
\makeatletter
\renewcommand{\ALG@name}{Giải thuật}
\renewcommand{\listalgorithmname}{Danh sách giải thuật}
\makeatother

\usepackage{listings}
\usepackage{xcolor}
\lstset{
  language=C++,
  basicstyle=\ttfamily\footnotesize,
  keywordstyle=\color{blue}\bfseries,
  stringstyle=\color{red},
  commentstyle=\color{green!50!black}\itshape,
  breaklines=true,
  frame=single,
  numbers=left,
  numberstyle=\tiny\color{gray},
  tabsize=2
}

%----------------------------------------------------------------------------------------
%	7. FONT CHUNG & KHOẢNG CÁCH ĐOẠN
%----------------------------------------------------------------------------------------
\makeatletter
\renewcommand\normalsize{%
  \@setfontsize\normalsize{13pt}{17.55pt}%
  \abovedisplayskip 6pt plus 2pt minus 2pt
  \abovedisplayshortskip \z@ plus 2pt
  \belowdisplayshortskip 6pt plus 2pt minus 2pt
  \belowdisplayskip \abovedisplayskip
}
\makeatother
\normalsize
\setlength{\parskip}{3pt}
\setlength{\parindent}{15pt}
\usepackage{indentfirst}

\renewcommand{\cftfigpresnum}{Hình~}
\setlength{\cftfignumwidth}{4.0em}

\renewcommand{\cfttabpresnum}{Bảng~}
\setlength{\cfttabnumwidth}{4.0em}

\usepackage{thesis}
\usepackage{setspace}

\cftsetpnumwidth{2em}
\cftsetrmarg{3em}
\setlength{\cftsecnumwidth}{3em}
\setlength{\cftsubsecnumwidth}{4em}

\setstretch{1.15}

%----------------------------------------------------------------------------------------
%	8. TÀI LIỆU THAM KHẢO
%----------------------------------------------------------------------------------------
%  "Tài liệu tham khảo" — không hiện header phải
\fancypagestyle{reflist}{%
  \fancyhf{}
  \fancyhead[L]{\fontsize{12pt}{14pt}\itshape\selectfont \makebox[\textwidth][l]{\datn}}
  \fancyfoot[L]{\fontsize{12pt}{14pt}\itshape\selectfont \nhomten}
  \fancyfoot[R]{\fontsize{12pt}{14pt}\selectfont \thepage}
  \renewcommand{\footrulewidth}{0.4pt}
}


\makeatletter
\patchcmd{\thebibliography}
  {\@mkboth{\MakeUppercase\bibname}{\MakeUppercase\bibname}}
  {\@mkboth{\MakeTextUppercase{\bibname}}{\MakeTextUppercase{\bibname}}\thispagestyle{reflist}}
  {}{}
\makeatother